\subsection{The Regime Prerequisite}

Our results establish that sampling-based consistency methods carry an implicit regime prerequisite: the model must operate in a stochastic hallucination regime, where wrong answers arise from generative variance rather than from confident but incorrect beliefs. When this prerequisite is violated --- as it is for Llama-3-8B-Instruct on structured factual QA --- the consistency signal collapses to noise.

This is not a failure of the consistency idea per se; it is a failure of the regime assumption. The original SMC papers \cite{Kuhn2023Semantic,Manakul2023SelfCheckGPT} validate on base language models performing open-ended generation, where the stochastic regime is a reasonable default. The implicit assumption is that a model generating ``Paris'' vs.\ ``Berlin'' vs.\ ``Rome'' across samples for the same question is exhibiting hallucination-correlated variance. This assumption breaks when RLHF training \cite{Ouyang2022Training} conditions the model to produce confident, consistent outputs regardless of factual accuracy.

\subsection{Why RLHF Induces Systematic Confabulation}

RLHF optimizes for human preference, which rewards fluent, confident, and helpful responses. A model that hedges or produces varied outputs may be rated less positively than one that commits to a specific answer. The training pressure therefore suppresses variance in outputs --- not selectively for correct answers, but uniformly. The result is a model that generates the same wrong answer consistently across samples, which looks identical to a model that generates the same right answer consistently.

This is the systematic confabulation regime: the model's outputs are internally consistent but not reliably factual. The consistency signal, which in the stochastic regime tracks uncertainty, now tracks only the model's trained tendency toward confident output --- a property that is orthogonal to factual accuracy.

\subsection{Implications for Hallucination Detection}

Our findings have several implications for practitioners deploying hallucination detectors on instruction-tuned models.

First, \emph{SMC methods should not be assumed to work on RLHF-fine-tuned models without empirical validation}. The regime assumption should be checked --- for example, by measuring mean consistency scores across labels before deploying the detector. A mean difference smaller than the within-group standard deviation is a warning sign.

Second, \emph{alternative signals may be needed for the systematic confabulation regime}. Consistency-based signals fail precisely because the model is consistent. Alternative approaches might include: retrieval-augmented verification (checking outputs against external knowledge), probing-based methods (examining internal representations rather than output distributions), or calibration-based methods that assess confidence without relying on output variance \cite{Kadavath2022Language,Xiong2023Can}.

Third, \emph{benchmark design matters}. HaluEval's structured factual QA format --- short factual questions with binary correct/hallucinated labels --- may particularly favor the systematic confabulation regime. Longer-form generation tasks, or tasks requiring multi-step reasoning where errors accumulate stochastically, may exhibit more variance. We encourage future work to evaluate SMC methods across a range of task types and model families.

\subsection{Limitations}

Our study has several limitations. We evaluate a single model (Llama-3-8B-Instruct) on a single benchmark (HaluEval QA). It is possible that other instruction-tuned models, other RLHF training recipes, or other task types would show different behavior. We do not compare against a base version of the same model (Llama-3-8B), which would provide direct evidence of the RLHF-induced regime shift. This comparison is left for future work.

Additionally, our analysis is observational: we infer the regime from distributional statistics rather than directly measuring model internals. A more direct test would involve probing the model's internal representations for factual uncertainty \cite{Kadavath2022Language}, which would require white-box access. Our evaluation is designed for the black-box setting, consistent with the original SMC papers.

Finally, the HaluEval dataset may have annotation noise --- some questions labeled as hallucinated may be borderline, reducing the signal available to any detector. However, since both SMC metrics fail by a large margin (AUROC near $0.49$ rather than, say, $0.48$), annotation noise is unlikely to be the primary explanation.

\subsection{When Might SMC Still Work?}

Our results do not imply that SMC methods are universally ineffective. We identify several conditions under which they may remain useful:

\begin{enumerate}
  \item \textbf{Base language models}: Models without RLHF fine-tuning may retain the stochastic hallucination regime, especially on open-ended generation tasks.
  \item \textbf{Long-form generation}: Tasks requiring extended reasoning chains may accumulate stochastic errors even in instruction-tuned models, reintroducing variance.
  \item \textbf{Low-resource or out-of-distribution domains}: When the model lacks strong priors about the correct answer, even RLHF-tuned models may produce varied outputs.
  \item \textbf{Calibrated models}: If RLHF training can be modified to preserve calibration \cite{Xiong2023Can}, the stochastic regime may be recoverable.
\end{enumerate}

We release our SMC infrastructure to support evaluation of these hypotheses.
